\chapter{系统设计}

\section{系统架构设计}

结合系统功能需求与技术栈特点，本系统设计为标准化、可扩展的分层架构，自下而上分为数据采集层、传输层、计算层、数据存储层以及应用展示层，明确各模块的职责与交互逻辑。

\subsection{技术架构}

系统采用前后端分离架构，后端基于 Spring Boot 2.7.15 + Spring Data JPA 构建 RESTful API 服务，前端基于 Vue 3 Composition API + Vite 8 + TypeScript + Tailwind CSS v4 + ECharts 5.6 构建响应式可视化界面。数据库采用 MariaDB，缓存层使用 Redis（Lettuce 连接池）配合内存缓存构建二级缓存体系。

在技术选型过程中，各层组件均经过对比权衡：前端在 React 与 Vue 之间选择了团队更为熟悉的 Vue 生态；后端在 Spring Boot 与 Flask 之间选择了前者，因为其在数据持久化与 API 构建方面生态更成熟；数据库在 MySQL 与 MariaDB 之间选择了 MariaDB，其完全兼容 MySQL 语法，且在性能与存储引擎上具备优势。

\begin{figure}[H]
    \centering
    \includegraphics[width=0.6\textwidth]{E:/download/系统架构图.drawio.png}
    \caption{系统功能模块图}
    \label{fig:function-modules}
\end{figure}

\subsection{数据采集层}

基于 MediaCrawler 框架实现 B 站与 NGA 双平台的数据采集。B 站爬虫通过搜索关键词获取视频列表及评论数据，采集字段包括 video\_id、title、desc、comment content、like\_count、coin\_count、favorite\_count、share\_count 等。NGA 爬虫通过论坛搜索获取帖子及回复数据，采集字段包括 thread\_id、title、author、content、replies 等。采集频率可灵活配置，定时增量执行。

\subsection{数据传输层}

传输层通过 Kafka 消息队列实现采集端与处理端的异步解耦。爬虫采集的原始数据封装为 JSON 格式后推送至 Kafka 统一主题，Spark Streaming 以 Consumer 身份消费数据。Kafka 作为缓冲层解决了采集端与处理端速度不匹配的问题，避免了数据堆积与丢失。依赖 ZooKeeper 协调节点状态，保证高吞吐与稳定性。

\subsection{计算层}

计算层是系统的核心处理单元，基于 Spark 3.3.0 + Scala 2.12.15 实现流式处理，通过 Kafka 消费原始数据并按批次执行清洗、计算与分析任务。该层按功能划分为四大模块，各模块职责与设计如下。

\subsubsection{数据清洗与标准化模块}

该模块负责将 B 站与 NGA 双平台的异构原始数据转换为结构统一的标准化 DataFrame，是整个分析流程的基础。清洗流程按顺序包含以下环节：

\begin{enumerate}
    \item \textbf{消息解析与基础校验}：解析 Kafka JSON 消息为 DataFrame，剔除 platform、raw\_id、author、publish\_time 等关键字段为空的脏数据行。
    \item \textbf{文本清洗}：使用正则表达式去除 HTML 标签、URL 链接、特殊符号、表情符号与非中文数字字符，去除多余空白符，统一编码格式。空文本标记为 null 并在后续过滤。
    \item \textbf{时间标准化}：将各平台原始时间字符串统一转换为 \texttt{yyyy-MM-dd HH:mm:ss} 格式的 Timestamp 类型，解析失败的行被丢弃。
    \item \textbf{字段映射与默认值填充}：将 platform 映射为整型（bilibili=1, nga=0），所有互动数值字段在 null 或负数时默认补 0，布尔字段转为 0/1 整型。
    \item \textbf{文本长度过滤}：清洗后文本长度不足 5 个字符的记录被过滤，保证语料质量。
\end{enumerate}

清洗完成后，双平台数据在字段名称、数据类型与语义上实现完全对齐，为后续热度计算与情感分析提供标准化的输入。

\subsubsection{热度计算模块}

热度计算模块采用分平台差异化加权策略，分别计算原始热度值后进行 Min-Max 归一化处理，使双平台热度具备可比性。

B 站平台支持点赞、评论、投币、收藏、分享等多种互动行为，原始热度公式为：

\[
Hot_{\text{raw}} = W_1 \cdot L + W_2 \cdot C + W_3 \cdot (Co + F + S) + 1.0
\]

其中 $L$ 为点赞数，$C$ 为评论数，$Co$、$F$、$S$ 分别为投币数、收藏数与分享数，权重当前配置为 $W_1{=}1.0$、$W_2{=}1.5$、$W_3{=}1.2$，基础值 1.0 确保零互动内容具有基准热度。

NGA 论坛仅支持评论互动，热度公式简化为：

\[
Hot_{\text{raw}} = W_2 \cdot C + 1.0
\]

归一化阶段采用分平台独立 Min-Max 策略，消除 B 站（高互动量）与 NGA（低互动量）之间的数值基数差异。B 站归一化区间为 $[0, 50000]$，NGA 为 $[0, 5000]$，归一化后截断至 $[0, 1]$ 区间得到 \texttt{hot\_norm}，直接作为最终热度 \texttt{hot\_score} 存储。

\subsubsection{情感分析模块}

情感分析模块基于外部 StructBERT 预训练模型服务实现。该服务以 Python Flask 框架独立部署在 localhost:8001 端口，加载 ModelScope 平台的 \texttt{damo/nlp\_structbert\_sentiment-classification\_chinese-base} 模型。

在 Spark 清洗流程中，通过自定义 UDF（User Defined Function）逐条将清洗后的文本内容发送至情感服务进行 HTTP POST 请求，服务端返回情感标签（positive/negative/neutral）与置信度分数。系统将结果映射为 $[-1.0, 1.0]$ 区间的 \texttt{sentiment\_score}：正值表示正面情感，负值表示负面情感，0 表示中性。输入文本长度限制为 200 字符以控制推理延迟，超出部分截断处理。

情感得分与热度数据一并写入数据库，为后续舆情分析中的情感趋势追踪、口碑量化评估提供支撑。

\subsubsection{热度预测模块}

热度预测模块基于 XGBoost 分类模型实现，以独立 Python Flask 服务（端口 5000）形式部署，通过 25 维特征对未来 2 小时窗口的热度趋势进行分类预测。

特征向量分为四类：
\begin{enumerate}
    \item \textbf{窗口聚合特征}（6 维）：当前 2 小时窗口内的 total\_hot、total\_hot\_raw、count、avg\_like、avg\_comment、avg\_sentiment，以及 keyword\_enc 编码。
    \item \textbf{历史滑动窗口特征}（12 维）：前 1--3 个窗口的 total\_hot、avg\_like、avg\_comment、avg\_sentiment 滞后值，用于捕捉时序依赖关系。
    \item \textbf{统计特征}（3 维）：热度变化率 total\_hot\_change\_rate，最近 6 个窗口的滚动均值与标准差，用于刻画波动趋势。
    \item \textbf{时间特征}（4 维）：hour、day\_of\_week、is\_weekend、month，用于捕获周期性规律。
\end{enumerate}

模型输出三分类结果：上升（热度变化率 $> +10\%$）、平稳（$[-10\%, +10\%]$）、下降（$<-10\%$），并附带各类别概率。基于概率加权估算预测变化率，最终计算预测热度值 \texttt{predicted\_hot\_raw}。该预测结果在前端趋势图中以虚线方式展示，辅助运营人员提前感知舆情走势。

\subsection{数据存储层}

采用 MariaDB 作为主存储，表名为 \texttt{standardized\_data}，存储经过 Spark 清洗标准化后的舆情全量数据。该表包含 32 个字段，涵盖平台标识、文本内容、互动指标、热度指标、情感得分等，并通过索引优化查询效率。Redis 作为缓存层，存储 API 查询结果，默认 5 分钟过期。同时在应用层维护 5 秒过期的本地内存缓存，进一步减少重复数据库查询。

\subsection{应用展示层}

后端基于 Spring Boot 2.7.15 运行于 8080 端口，提供 6 个 RESTful API 端点，均位于 \texttt{/api/opinion/} 路径下。各端点功能如下：\texttt{/metrics} 提供平均热度指数、环比变化、情感占比与平台分布比例等核心指标；\texttt{/channel} 返回 B 站与 NGA 双平台的数据量占比；\texttt{/trend} 按 12 等分时间窗口聚合声量走势数据，末尾追加预测点；\texttt{/words} 基于 Ansj 分词器对文本进行中文分词与词频统计，返回权重归一化的热词列表；\texttt{/list} 按时间倒序与热度降序综合排序，返回最新 6 条舆情信息；\texttt{/health} 提供健康检查。所有接口接收 \texttt{keyword} 与 \texttt{timeRange} 两个查询参数，后端根据 timeRange 解析为具体的 Date 范围进行数据库查询，返回统一格式的 \texttt{ApiResponse<T>} 响应体（含 code、message、data 三个字段）。

前端开发服务器运行于 5173 端口，使用 Axios 库向后端发起 HTTP 请求。前端在关键词或时间范围切换时自动触发对应 API 调用，并通过 Vue 的响应式系统将数据传递至各子组件。可视化层面采用 ECharts 5.6 进行图表渲染，共使用三种图表类型：折线图用于展示声量走势与迷你趋势，支持平滑曲线与渐变面积填充；饼图用于展示渠道声量分布，采用环形样式并支持悬停高亮；词云图基于 echarts-wordcloud 插件实现，圆形布局，词频越高字号越大。所有图表均具备自适应缩放能力，绑定浏览器 resize 事件实时更新尺寸。

\section{数据库设计}

系统数据库基于 MariaDB 设计，遵循数据结构化、冗余最小化、查询高效化原则。完整数据库 ER 关系如图~\ref{fig:database-er} 所示。

\begin{figure}[H]
    \centering
    \includegraphics[width=0.2\textwidth]{E:/download/deepseek_mermaid_20260423_72a9a0.png}
    \caption{数据库ER图}
    \label{fig:database-er}
\end{figure}

\subsection{核心数据表结构}

系统核心数据表为 \texttt{standardized\_data}，存储经 Spark 预处理后的双平台统一结构数据。该表共 32 列，以 \texttt{raw\_id} 为物理主键，对 \texttt{keyword} 和 \texttt{publish\_time} 建立联合索引以优化时间范围查询，具体表结构设计如下。

\begin{figure}[H]
    \centering
    \includegraphics[width=0.3\textwidth]{E:/download/标准数据输出表.png}
    \caption{标准化数据表结构}
    \label{fig:standard-data-table}
\end{figure}

各主要字段说明：
\begin{itemize}
    \item \texttt{raw\_id}：原始数据唯一标识，物理主键
    \item \texttt{platform}：平台标识（1=B站，0=NGA）
    \item \texttt{keyword}：所属搜索关键词
    \item \texttt{title\_clean} / \texttt{content\_clean}：清洗后的标题与正文
    \item \texttt{publish\_time}：发布时间（\texttt{yyyy-MM-dd HH:mm:ss}）
    \item \texttt{like\_count}、\texttt{comment\_count}、\texttt{coin\_count}、\texttt{favorite\_count}、\texttt{share\_count}：互动指标
    \item \texttt{hot\_raw}：原始热度加权值
    \item \texttt{hot\_norm}：Min-Max 归一化后的热度值
    \item \texttt{hot\_score}：最终热度评分（当前与 \texttt{hot\_norm} 一致）
    \item \texttt{sentiment\_score}：情感得分（$[-1.0, 1.0]$）
    \item \texttt{text\_length}：文本长度
\end{itemize}

\section{可视化界面设计}

Web 界面遵循简洁、直观、高效的设计原则，面向游戏运营人员与舆情研究人员打造轻量化使用体验。系统采用浅色主题（背景色 \#f1f5f9），卡片式布局配合圆角与阴影效果，整体风格明亮清晰。

\begin{figure}[H]
    \centering
    \includegraphics[width=0.8\textwidth]{E:/展示图片1.png}
    \caption{Web 界面展示 1}
    \label{fig:web-interface-1}
\end{figure}

主页面布局分为四个区域：

\begin{enumerate}
    \item \textbf{顶部导航栏}：左侧提供关键词下拉选择器（21 个预设手机游戏关键词），右侧提供时间范围按钮组（1天至全部共 8 个选项）。导航栏采用毛玻璃效果的白色半透明背景，通过关键词与时间范围的组合筛选控制页面所有卡片的数据范围。关键词下拉菜单支持点击外部区域自动关闭；时间范围在桌面端以按钮组形式展示，移动端自动切换为下拉选择框。

    \begin{figure}[H]
        \centering
        \includegraphics[width=0.8\textwidth]{E:/顶部导航栏.png}
        \caption{顶部导航栏展示}
        \label{fig:top-nav}
    \end{figure}

    \item \textbf{核心指标区}：左右并列排放两张指标卡片，提供舆情数据的顶层量化概览。

    \textbf{平均热度指数卡片}展示当前关键词在选定时间范围内所有舆情内容的平均热度值、与上一同周期的环比变化百分比，以及迷你热度趋势折线图。用户通过该卡片可快速判断整体舆情热度水平及其走向——热度是处于上升通道还是下降通道。

    \textbf{情感分布卡片}以三色百分比进度条展示正面、中性、负面三类情感的占比分布，其中正面以绿色显示、中性以黄色显示、负面以红色显示。用户通过该卡片可一键掌握所选关键词下的舆论情感倾向，判断口碑是积极、消极还是以中性为主。

    \item \textbf{可视化图表区}：包含三张图表卡片，分别从渠道结构、时间趋势与文本焦点三个维度展示舆情数据。

    \textbf{渠道声量分布饼图}以环形饼图展示 B 站与 NGA 在指定时间范围内的数据量占比，悬停时可查看各平台的具体比例。用户可据此判断不同平台对特定关键词的讨论活跃度差异。

    \textbf{实时热度与预测折线图}以折线图展示时间范围内的声量走势，同时叠加 XGBoost 模型预测的下一个 2 小时窗口热度值（以虚线与预测点标记），右上角标注趋势方向（上升/平稳/下降）及置信度。用户可据此了解历史热度变化规律，并对近期走势形成预判。

    \textbf{热点词云图}以词云形式展示高频关键词，权重越高的词字号越大，悬停时显示词语及其权重。用户可直观识别当前讨论的核心话题与焦点词汇。

    \item \textbf{信息流区}：以排行榜形式展示最新舆情信息列表，按时间与热度综合排序，每条信息包含标题、来源平台、发布时间描述、情感标签与热度数值，最多展示 6 条。用户可通过该区域快速浏览最新讨论动态，了解当前最受关注的舆情内容。
\end{enumerate}

\begin{figure}[H]
    \centering
    \includegraphics[width=0.8\textwidth]{E:/展示图片2.png}
    \caption{Web 界面展示 2}
    \label{fig:web-interface-2}
\end{figure}

交互设计上，切换关键词或时间范围时触发所有卡片统一刷新，加载过程中各卡片独立展示旋转加载动画。页面整体采用响应式布局，桌面端图表区为两列排列，移动端自动堆叠为单列，适配不同屏幕尺寸。

\section{系统特点}

从整体架构来看，前后端分离使各层可独立开发与部署，系统具备良好的可扩展性。分层架构设计使数据采集、处理、存储与展示职责清晰，便于后续添加新的数据源或分析模块。

功能层面，系统覆盖了从数据采集到可视化展示的完整链路：双平台多渠道数据采集保证信息全面性；基于 Structured Streaming 的流式处理确保数据时效性；XGBoost 模型预测提供前瞻性舆情洞察；多维度的可视化图表使数据一目了然。

性能方面，系统采用 Redis + 内存缓存二级缓存体系，有效降低数据库查询压力。数据库对 keyword 与 publish\_time 字段建立联合索引，优化时间范围查询效率。后端接口响应时间在缓存命中场景下可控制在毫秒级。
